\section{The value of disclosure for safe trading}\label{sec:disclosure}

The companion paper measures the value of disclosure in samples needed to identify the cohorts \citep{niroula2026}. For a pre-trade control, the relevant value is the safe trading capacity it unlocks.

\begin{theorem}[Decision value of disclosure]\label{thm:disclosure}
Under the risk contract of Section~\ref{sec:contracts}, compare an observation contract $O$ with a richer contract $O'$ that reveals everything $O$ reveals and more, both with exact data.
\begin{enumerate}[label=(\roman*)]
\item \emph{Monotonicity.} The compatible sets are nested, $\Theta_{O'}\subseteq\Theta_O$, so $\Phi^+_{O'}\le\Phi^+_O$ and the admitted level of Theorem~\ref{thm:robustguard} satisfies $\alpha^*_{O'}\ge\alpha^*_O$.
\item \emph{Strictness.} $\alpha^*_{O'}>\alpha^*_O$ if and only if $\alpha^*_O<1$, $c_I>0$, and $O'$ excludes every compatible economy with $\Phi=\Phi^+_O$. Under the window contract that economy is the upper principal representation of Theorem~\ref{thm:window}: disclosure has decision value exactly when it rules out a cohort at the memory bound with the extremal weight.
\item \emph{Long memory.} Under the window contract with memory bound $\bar r$, $\alpha^*_O\to0$ as $\bar r\to1$ for any fixed window (Theorem~\ref{thm:window}(iii)), whereas audited exposure, which reveals $\Phi$, admits the oracle capacity $\alpha^\circ$. The decision value of audited disclosure therefore tends to the whole oracle capacity as memory lengthens.
\item \emph{Net value.} With value $w$ per unit of admitted capacity and disclosure cost $c(O')$, the net value $w(\alpha^*_{O'}-\alpha^*_O)-c(O')$ is zero or negative whenever $O$ already identifies $\Phi$ (a window whose moment vector lies on the boundary of the moment space) or the budget does not bind.
\end{enumerate}
With noisy data the confidence sets of two contracts need not be nested; intersecting two certificates, each at level $1-\gamma/2$, restores monotonicity at level $1-\gamma$.
\end{theorem}

The theorem is set inclusion plus the closed form of Theorem~\ref{thm:robustguard}; its content is the location of the value. Disclosure is worth most where public certificates are weakest, in markets whose participants have long memories. In simulation (test T07; Table~\ref{tab:discsim}), with two classes and one of them slow, audited exposure adds 0.027 of admitted capacity to a public certificate when the slow class's memory rate is 0.6 and 0.122 when it is 0.995, out of an oracle capacity of 0.667. Class-tagged flows measured with 10\% noise add 0.02--0.03 at memory rates up to 0.97 but \emph{lower} capacity at 0.995, because a class's exposure $q=w/(1-r)$ is badly conditioned when $r$ is near one: disclosure of a noisy proxy can be worth less than nothing. The companion paper's evidence that Binance's and the CFTC's class positions pin down memories that prices leave unrestricted concerns class-level positions, which are proxies, not records of individual algorithms.

\begin{table}[t]\centering\small
\begin{threeparttable}
\caption{Admitted capacity under three observation contracts (simulation)}\label{tab:discsim}
\begin{tabular}{lcccc}
\toprule
Memory rate of the slow class & 0.60 & 0.90 & 0.97 & 0.995\\
\midrule
Public responses (certificate, Theorem~\ref{thm:certcost}) & 0.640 & 0.623 & 0.601 & 0.544\\
Class-tagged flows, 10\% noise & 0.662 & 0.655 & 0.631 & 0.507\\
Audited exposure (oracle) & 0.667 & 0.667 & 0.667 & 0.667\\
\midrule
Value of audited exposure over public & 0.027 & 0.044 & 0.066 & 0.122\\
\bottomrule
\end{tabular}
\begin{tablenotes}\footnotesize\item Two classes ($q=0.06$ with the stated memory rate, $q=0.12$ with rate $0.5$), $a=0.4$, $\kappa=1$, position 1, shock bound $0.48$, budget 1; public: 500 reset experiments, per-tap noise $0.05$, certificate at 95\% with the true memory bound; tagged: 2{,}000 periods of each class's flow with 10\% measurement noise, two-stage least squares as in the companion paper, 95\% delta-method bound; 300 replications. Entries are mean admitted fractions $\alpha^*$. Kind: recomputed (test T07).\end{tablenotes}
\end{threeparttable}
\end{table}

\paragraph{Pricing versus guarding.} A conservative bound is the right input for a guard and the wrong input for a price. The corrective rebate of Proposition~\ref{prop:equil} depends on $S^2$; computed from an upper bound that overstates $S$ by 10\% it over-subsidizes reversion and loses 2.4--2.8\% of first-best welfare, and by 50\%, 70--80\% (test T12). Guards should use certified upper bounds; rebates should use point estimates.
